\documentclass[11pt,a4paper]{article}

\usepackage[main=spanish]{babel}
\usepackage[utf8]{inputenc}
\usepackage[T1]{fontenc}

\usepackage{amsmath, amssymb, amsfonts}
\usepackage{mathtools}
\usepackage{geometry}
\usepackage{graphicx}
\usepackage{setspace}
\usepackage{hyperref}
\usepackage{graphicx}
\usepackage{appendix}
\usepackage{listings}
\usepackage{xcolor}

\geometry{margin=2.5cm}
\onehalfspacing
\title{Laboratorio 2: Lagrange y KKT}
\author{Adrián Janus Gonzalez Adamuz}
\date{\today}
\lstset{
  inputencoding=utf8,
  extendedchars=true,
  basicstyle=\ttfamily\small,
  numbers=left,
  numberstyle=\tiny,
  frame=single,
  breaklines=true,
  showstringspaces=false,
  tabsize=2,
  literate=
    {á}{{\'a}}1 {é}{{\'e}}1 {í}{{\'\i}}1 {ó}{{\'o}}1 {ú}{{\'u}}1
    {Á}{{\'A}}1 {É}{{\'E}}1 {Í}{{\'I}}1 {Ó}{{\'O}}1 {Ú}{{\'U}}1
    {ñ}{{\~n}}1 {Ñ}{{\~N}}1
  }	

\begin{document}
\maketitle

\section{Lagrange}
\subsection{Ejemplo 1}
Consideramos el problema de maximizar $f(x,y) = 2x + 3y$ sujeta a la restricción $\sqrt{x} + \sqrt{y} = 5$ 
\begin{itemize}
\item Utiliza la computadora para graficar la función y la restricción. 
\item Intenta usar los multiplicadores de Lagrange para resolver el problema. 
\item Usa la gráfica de la parte (a) para intentar encontrar el punto maximizador deseado. 
\item ¿Cuál es la significancia del punto $f(9,4)$?
\item Explica por qué falla el método de los multiplicadores de Lagrange.

\end{itemize}

\subsubsection{Grafica y restricción}
\begin{figure}[hbtp]
\centering
\includegraphics[scale=0.5]{../../Imágenes/Sin título.png}
\caption{Grafica con restricción}
\end{figure}
De aqui se ve claramanete como el máximo de la función se encuentra en el (0,25) y hay otro candidato en el (25,0). 
Evaluando en la función
\[
f(0,25) =2(0) + 3(25)=75 \quad f(25,0)=2(25) +3(0) = 50,
\]

\subsubsection{Desarollo Lagrange}
Tenemos que $g(x,y) = \sqrt{x} + \sqrt{y} - 5$ como la restricción, ahora sacamos el gradiente de ambas fucniones:
\[
\nabla f(x,y) = \begin{pmatrix}
2 & 3
\end{pmatrix}
\quad
\nabla g(x,y) = \begin{pmatrix}
\frac{1}{2\sqrt{x}}
&
\frac{1}{2\sqrt{y}}
\end{pmatrix}
\]
Recordamos la condición de lagrange
$\nabla f -\lambda \nabla g = 0$
Entonces tenemos el sistema
\[
2 = \lambda \frac{1}{2\sqrt{x}}, \quad 3= \lambda \frac{1}{2 \sqrt{y}}
\]
despejando 
\[
\lambda = 4 \sqrt{x} \quad \lambda = 6 \sqrt{y}
\]
igualando las lambdas tenemos que
\[
\sqrt{x} = \frac{3}{2} \sqrt{y}
\]
Si sustituimos esta igualdad en la restriccion $g(x,y)$
\[
\frac{3}{2} \sqrt{y} + \sqrt{y} = 5 \Rightarrow \frac{5}{2} \sqrt{y} = 5 \Rightarrow \sqrt{y} = 2 \Rightarrow y = 4
\]
Entonces $\sqrt{x} = 3 \Rightarrow x=9$
Por lo tanto el punto que nos da Lagrange es $(x,y) = (9,4)$
Evalando en la función
\[
f(9,4) = 2(9) +3(4) = 18 + 12 = 30
\]
\subsubsection{Importancia de $f(9,4)$}
El punto $(9,4)$ es significativo por dos razones:
\begin{enumerate}
\item Es factible porque satisface la restricción $\sqrt{9} + \sqrt{4} = 3+2 = 5$
\item Es el punto crítico interior que produce el método de lagrange
\end{enumerate}
\subsubsection{¿Por qué falla Lagrange?}
El problema no es que Lagrange no sirva, sino que no se cumplen las hipótesis necesarias.
Lagrange solo sirve si
\begin{enumerate}
\item El punto óptimo sea interorior respecto a la frontera
\item la restricción se diferenciable y $\nabla g \neq 0$ en el punto.
\end{enumerate}
\subsection{Ejemplo 2}
Consideramos el problema de minimizar la función $f(x,y)= x$ en la curva $y^2 + x^4 - x^4 =0$ denominada curva piroforme
\begin{enumerate}
\item Utiliza la computadora para graficar la función y la restricción
\item ntenta usar los multiplicadores de Lagrange para resolver el problema
\item Muestra que el mínimo de la función es $f(0, 0) = 0$ pero la condición de Lagrange$\nabla f(0,0) = \lambda \nabla g(0,0)$ no se cumple para ningún valor de $\lambda$.
\item Explica porque falla el método de Lagrange. 
\end{enumerate}
\subsubsection{Grafica}
\begin{figure}[hbtp]
\centering
\includegraphics[scale=0.5]{../../Imágenes/grafica2.png}
\caption{Curva piroforme}
\end{figure}
\subsection{Multiplicadores de Lagrange}
\[
\nabla f(x,y) = \lambda \nabla g(x,y) \Rightarrow \nabla f(x,y) = (1,0), \quad \nabla g(x,y) = (4x^3 - 3x^2, 2y)
\]
La condición de lagrange queda 
\[
(1,0) = \lambda (4x^3 - 3x^2, 2y)
\]
De la segunda ordenada vemos que que $2y\lambda = 0$ de manera rápida vemos que si $\lambda = 0$ lleva a que $0=1$ entonces si sustituimo $y=0$ en la restricción 
\[
0 + x^4 - x^3 = x^3(x-1) =0 \Rightarrow x=0 \quad o \quad x = 1
\]
Entonces los únicos puntos factibles son $y=0$ son $(0,0)$ y $(1,0)$
Evaluamos en $\nabla g$ 
\[
\nabla g(1,0) =(1,0), \quad g(0,0) = (0,0)
\]
3.De manera gráfica y analítica el mínimo se encuentra en el (0,0) $f(0,0) = 0$
pero $\nabla f(0,0) = (1,0)$ y $\nabla g(0,0) = (0,0)$ es decir tendriamos la condición de Lagrange siguiente
$(1,0) = \lambda(0,0) !$ Por lo tanto
\[
\nexists \lambda : \nabla f(0,0) = \lambda \nabla g(0,0)
\]
\subsubsection{¿Por qué falla Lagrange?}
Porque el método de multiplicadores de Lagrange requiere regularidad de la restricción en el punto óptimo: típicamente que 
g sea diferenciable y que $\nabla g(x^*, y^*) \neq (0,0)$
En $(0,0)$ ocurre exactamente lo opuesto $\nabla g(0,0) = (0,0)$ y geometricamente la curva tiene un "pico" en el $(0,0)$
haciendola no diferenciable en ese punto.

\section{KKT}
\subsection{Ejemplo 1}
Encontramos el máximo de $f(x,y) = 3x + 4y$ sujeto a las restricciones \[g_1(x,y) = x^2 + y^2 \leq 4\]
\[ g_2(x,y) = x \geq 1 \]
\begin{enumerate}
\item Enumera las condiciones de Kuhn-Tucker para este problema.
\item Resuelve las ecueaciones anteriores
\item Grafica la función objetivo
\end{enumerate}
Planteamos el Lagrangiano para maximización con restricciones $\leq$ 0:
\[
\mathcal{L}(x,y,\lambda_1, \lambda_2) = 3x + 4y - \lambda_1(x^2 + y^2 - 4)-\lambda_2 (1-x) 
\] con $\lambda_1, \lambda_2 \geq 0$.
Las condiciones del KKT son:
\newline
Factibilidad primal
\[x^2 + y^2 \leq 0, \quad 1-x\leq 0\]
Factibilidad dual
\[\lambda_1 \geq 0, \lambda_2\geq 0 \]
Estacionariedad $\nabla_{x,y} \mathcal{L} = 0$
\[\frac{\partial \mathcal{L}}{\partial x} = 3-2\lambda_1x+\lambda_2 = 0\]
\[\frac{\partial \mathcal{L}}{\partial y} = 4 -2\lambda_1y = 0.\]
Complementariedad
\[\lambda_1(x^2+y^2 -4)=0\]
\[\lambda_2(1-x)=0\]
\subsubsection{Resolución del sistema}
Tomamos $g_1(x,y)=x^2+y^2-4\le 0$ y $g_2(x,y)=1-x\le 0$. 
\[
\mathcal L(x,y,\lambda_1,\lambda_2)=3x+4y-\lambda_1(x^2+y^2-4)-\lambda_2(1-x),
\qquad \lambda_1,\lambda_2\ge 0.
\]
Resolviendo el sistema:
\[
\frac{\partial \mathcal L}{\partial x}=3-2\lambda_1 x+\lambda_2=0,\qquad
\frac{\partial \mathcal L}{\partial y}=4-2\lambda_1 y=0.
\]
De $4-2\lambda_1 y=0$ se obtiene $\lambda_1 y=2$, luego $\lambda_1>0$ y por complementariedad $x^2+y^2=4$.

\emph{Caso 1:} $x>1 \Rightarrow \lambda_2=0$. Entonces
\[
3-2\lambda_1 x=0\Rightarrow x=\frac{3}{2\lambda_1},\qquad
y=\frac{2}{\lambda_1}.
\]
Imponiendo $x^2+y^2=4$:
\[
\left(\frac{3}{2\lambda_1}\right)^2+\left(\frac{2}{\lambda_1}\right)^2=4
\ \Rightarrow\ 
\lambda_1=\frac54,
\]
y por tanto
\[
(x^*,y^*)=\left(\frac{6}{5},\frac{8}{5}\right),\qquad \lambda_2=0.
\]

\emph{Caso 2:} $x=1$. De $x^2+y^2=4$ se sigue $y=\pm\sqrt3$; pero $y=-\sqrt3$ da $\lambda_1<0$ (imposible). Con $y=\sqrt3$ se tiene $\lambda_1=\frac{2}{\sqrt3}$ y
\[
3-2\lambda_1(1)+\lambda_2=0 \Rightarrow \lambda_2=2\lambda_1-3=\frac{4}{\sqrt3}-3<0,
\]
contradicción. Por tanto no hay solución KKT con $x=1$.

Concluimos que el máximo ocurre en
\[
(x^*,y^*)=\left(\frac{6}{5},\frac{8}{5}\right),\qquad f(x^*,y^*)=3\cdot\frac65+4\cdot\frac85=10,
\]
con multiplicadores $\lambda_1=\frac54$ y $\lambda_2=0$.

\newpage
\subsubsection{Gráfica}
\begin{figure}[hbtp]
\centering
\includegraphics[scale=1]{../../Imágenes/Graficas/3.png}
\end{figure}

\subsection{Ejemplo 2}
Encuentra el máximo la función $f(x,y) = 2x^2+ 3xy$ sujeta a las restricciones 
\[
g_1(x,y) = \frac{1}{2}x^2 + y \leq 4
\]
\[
g_2(x,y) =  y \geq 2
\]

Planteamos el Lagrangiano para maximización con restricciones $\leq 0$:
\[
\mathcal{L}(x,y,\lambda_1, \lambda_2)
=
2x^2+3xy
-\lambda_1\!\left(\frac{1}{2}x^2 + y-4\right)
-\lambda_2 (2-y)
\]
con $\lambda_1, \lambda_2 \geq 0$.

Las condiciones del KKT son:
\newline
Factibilidad primal
\[
\frac{1}{2}x^2 + y-4 \leq 0, \quad 2-y\leq 0
\]
Factibilidad dual
\[
\lambda_1 \geq 0, \lambda_2\geq 0
\]
Estacionariedad $\nabla_{x,y} \mathcal{L} = 0$
\[
\frac{\partial \mathcal{L}}{\partial x}
=
4x+3y-\lambda_1 x = 0
\]
\[
\frac{\partial \mathcal{L}}{\partial y}
=
3x-\lambda_1+\lambda_2 = 0
\]
Complementariedad
\[
\lambda_1\!\left(\frac{1}{2}x^2+y-4\right)=0
\]
\[
\lambda_2(2-y)=0
\]

\subsubsection{Resolución del sistema}

Como $y\ge 2$, de la ecuación de estacionariedad en $x$ se obtiene
\[
x(4-\lambda_1)+3y=0,
\]
lo que implica $\lambda_1>0$. Por complementariedad,
\[
\frac{1}{2}x^2+y=4.
\]

Tomamos el caso $y>2$, de donde $\lambda_2=0$. El sistema queda
\[
\begin{cases}
4x+3y-\lambda_1 x=0,\\
3x-\lambda_1=0.
\end{cases}
\]

De la segunda ecuación se obtiene $\lambda_1=3x$. Sustituyendo en la primera:
\[
4x+3y-3x^2=0.
\]

Usando la restricción activa:
\[
y=4-\frac{1}{2}x^2,
\]
se obtiene
\[
4x+3\left(4-\frac{1}{2}x^2\right)-3x^2=0,
\]
\[
\frac{9}{2}x^2-4x-12=0.
\]

Resolviendo,
\[
x=\frac{4\pm\sqrt{232}}{9}.
\]
Tomamos la raíz positiva,
\[
x^*=\frac{4+\sqrt{232}}{9},
\qquad
y^*=4-\frac{1}{2}(x^*)^2.
\]

Por lo tanto, el máximo se alcanza en
\[
(x^*,y^*)=\left(\frac{4+\sqrt{232}}{9},\;4-\frac{1}{2}\left(\frac{4+\sqrt{232}}{9}\right)^2\right).
\]
\subsection{Gráfica}
\includegraphics[scale=1]{../../Imágenes/Graficas/4.png} 
\end{document}